\section{标签、尺寸与正文内嵌}
\begin{example}{示例 31：标签位置}
\gradbar[label=outside]{64}

\gradbar[label=inside,
  height=2.5ex,
  right color=white]{64}

\gradbar[label=none]{64}
\end{example}

外部标签采用右对齐的槽位，默认在文字过长时扩展；同列应统一设置足够大的槽宽。
内部标签位于整条轨道右端附近，不跟随彩色条的末端移动。
\key{label=none} 则只显示轨道和填充。

\begin{example}{示例 32：单位、精度和字体}
\gradbar[max=5000,
  width=18mm,label width=6em,
  precision=0,unit={\,ms}]{1250}

\gradbar[width=18mm,
  precision=2,
  font={\small\bfseries}]{86.537}
\end{example}

\begin{example}{示例 33：宽度与高度}
\gradbar[width=18mm,
  height=1ex]{70}

\gradbar[width=26mm,
  height=2ex]{70}

\gradbar[width=26mm,
  label gap=1em,
  label width=3em]{70}
\end{example}

\key{width} 只指定轨道宽度。带外部标签时，总宽度约为轨道宽度、
\key{label gap} 和 \key{label width} 三者之和。
这些长度支持 mm、pt、em、ex 等单位。

\begin{example}{示例 33a：左侧标签与基线}
\gradbar[label=left,unit={\%},baseline=text]{92.7}
\qquad
\gradbar[label=left,baseline=bottom]{NA}
\end{example}

\key{label=left} 把格式化后的标签放在轨道左侧，并在 \key{label width} 槽位内右对齐。
它适合需要“数值在前、图形在后”的紧凑表格列；精度、单位、字体和缺失值规则保持不变。
\key{baseline=text} 按当前标签字体的度量对齐，\key{baseline=center} 和
\key{baseline=bottom} 分别按图形中心或轨道下边缘对齐；正的 \key{baseline shift}
会把图形向上移动。

\begin{example}{示例 34：直接放入正文}
Completion:
\gradbar[width=15mm,
  height=1.2ex,
  label width=3em,
  unit={\%},precision=0]{64}

The task is progressing.
\end{example}

\note{标签默认沿用周围文字的字体。普通表格行高仍由表格本身控制；
需要更宽松的行距时设置 \key{\arraystretch}。宏包不会修改整张表格的行距。}
